\section{Discussion}\label{sec:discussion}

\subsection{The transferable result is an operator distinction, not a cadence threshold}

The exponential frozen-assessment and forward-Euler controller updates, and the continuous-delay comparator, yield different spectral transitions on the same illustrative ecological plant. This is the central mathematical result. It cautions against transferring a stability boundary from one representation to another without matching (i) what is observed and when, (ii) which state estimate enters the rule, (iii) when the command changes, (iv) whether it is held between reviews, and (v) how command gain scales with the interval. In a management-strategy evaluation, those items should be explicit alongside operating-model uncertainty and implementation error (Punt et al., 2016).

The comparison does not show that real fisheries are unstable under annual reviews or that a longer interval is generally safer. The baseline exponential-update crossing is parameter-sensitive, the controller is uncalibrated, and the exponential protective result is only a scan over a finite interval grid. A fitted institutional threshold would require independent estimates of ecological response, observation and assessment error, decision and deployment lags, and the effective rule gain. Without them, the reported years are reproducible properties of a stated model, not policy prescriptions.

\subsection{Empirical diagnostics constrain stronger claims}

The baseline forage-fish screen contains no multiplicity-adjusted target-band detection, and the sensitivity variants do not change that decision. This is a useful negative diagnostic but not a test of a causal governance mechanism. Lomb--Scargle band power can be elevated by external forcing and structural change; an AR(1) null does not model all such alternatives. Nor does a spectrum alone reveal whether management responded to a signal before or after the ecological change.

The corrected reading of the review-interval test is particularly important. There is no observed interval spanning the modeled response windows, so there is no identifiable window-approach gradient. What is reproduced is an opposite-direction contrast between two proxy groups: seven regionally clustered 0.5-year assignments have larger null-excess than 35 annual assignments. Since the shorter proxy is farther from both modeled windows, this contrast does not support the proposed approach mechanism. The two-value coding, group-level rather than stock-level verification, and regional concentration leave confounding unresolved. ENSO-related quasi-periodicity is a candidate explanation, not a demonstrated cause. This distinction replaces any earlier shorthand suggesting that the cross-sectional test found a predicted gradient.

The 32-system search likewise has a narrow interpretation. No inventory candidate satisfies all four evidential requirements, but the inventory is a curated search rather than an exhaustive review, and missing station-level data could exclude otherwise relevant cases. The result is that a qualifying demonstration was not found in this search, not that institutional feedback cannot occur.

\subsection{A descriptive cod partition is not a causal decomposition}

Northern cod provides an instructive record because the collapse and subsequent low-stock period are well documented. However, a two-window plot cannot identify the causal contribution of harvest rules, predation, climate, recruitment, observation, or assessment structure. The strong-Allee surplus model is a local mathematical lens; unless fitted against alternatives and validated out of sample, it does not establish an Allee mechanism. Likewise, the constrained-mortality formulation discussed in prior notes cannot support numerical conclusions here until its source table, uncertainty, and complete computation are recovered. The manuscript therefore retains only the reproducible descriptive partition and lists the unreproduced mortality claims as open hypotheses in Supplementary S13.2.

\subsection{Requirements for a governance-feedback test}

A stronger empirical design should treat institutions as measured mechanisms rather than inferred labels. At minimum, it would require:
\begin{enumerate}
\item dated observations, assessment releases, rule decisions, implementation, and ecological response at a resolution finer than the candidate response window;
\item stock-specific evidence about the operative rule, compliance, command persistence, and the timing of responsive revisions;
\item a state-space model that separates observation error, process noise, external forcing, regime shifts, and structural assessment changes;
\item pre-specified comparisons among no-feedback, delayed-feedback, sample-and-hold, and ecological-only models, with held-out predictive scoring or a credible intervention design; and
\item uncertainty propagation for the estimated review map, rather than reporting only a single point crossing.
\end{enumerate}
These requirements are compatible with management-strategy evaluation and adaptive management practice (Walters, 1986; Punt et al., 2016), but are more demanding than detecting a spectral peak. Supplementary S7 sets out five prospective designs with estimands and decision rules; none is reported as executed.

\subsection{Limitations}

This edition has six important limitations. First, the focal numerical parameters are illustrative, with no joint field calibration. Second, local eigenvalue stability does not imply global nonlinear stability or acceptable transient performance. Third, dense-grid non-crossing results are numerical evidence only. Fourth, the RAM screen and \(\Tr\)-proxy assignments inherit database coverage and analyst coding; stock-specific primary verification is incomplete. Fifth, the stage-map reconstruction does not reproduce all archived trajectory windows, which remain provisional. Sixth, the case inventory and cod analysis do not identify a causal governance effect. Supplementary materials retain complete outputs and provenance so that these limitations are not hidden by a concise main text.

\section{Conclusion}\label{sec:conclusion}

Review cadence is a dynamical design variable when decisions are sampled and commands are held, but a cadence threshold is meaningful only for a fully specified and empirically supported update map. In the illustrative system, exponential frozen-assessment, Euler, and continuous-delay operators produce different crossings, while small joint parameter perturbations move or remove them. The analytic invariance and rapid-review results do not supply a stability guarantee at finite review intervals. The empirical screens do not reveal a window-approach gradient; the reproduced two-group association points in the opposite direction and is confounded by coarse, regionally clustered proxies. Northern cod remains a descriptive case. These findings motivate operator-aware, data-rich evaluations rather than a general claim that any particular governance interval stabilizes real resources.

\section*{References}\label{references}

Abaee, A. 2026. Delay-induced regime change in harvested stocks: the
mobilising and protective channels of institutional feedback, and the
review interval as control. Zenodo.
https://doi.org/10.5281/zenodo.22554217.

Adamson, M. W., and Hilker, F. M. 2020. Resource-harvester cycles caused by delayed knowledge of the harvested population state can be dampened by harvester forecasting. Theoretical Ecology, 13: 425--434.

{\AA}str{"o}m, K. J., and Wittenmark, B. 1997. Computer-Controlled Systems: Theory and Design, 3rd edn. Prentice Hall, Upper Saddle River, NJ.

Alkire, S., and Foster, J. 2011. Counting and multidimensional poverty
measurement. Journal of Public Economics, 95: 476--487.

Ashwin, P., Wieczorek, S., Vitolo, R., and Cox, P. 2012. Tipping points
in open systems: bifurcation, noise-induced and rate-dependent examples
in the climate system. Philosophical Transactions of the Royal Society
A, 370: 1166--1184.


Benjamini, Y., and Hochberg, Y. 1995. Controlling the false discovery
rate: a practical and powerful approach to multiple testing. Journal of
the Royal Statistical Society B, 57: 289--300.

Benjamini, Y., and
Yekutieli, D. 2001. The control of the false discovery rate in multiple
testing under dependency. Journal of the Royal Statistical Society B,
63: 425--441.

Branch, T. A. 2001. A review of orange roughy Hoplostethus atlanticus
fisheries, estimation methods, biology and stock structure. South
African Journal of Marine Science, 23: 181--203.

Brown, C. J., Fulton, E. A., Possingham, H. P., and Richardson, A. J. 2012. How long can fisheries management delay action in response to ecosystem and climate change? Ecological Applications, 22: 298--310.

Butterworth, D. S. 2007. Why a management procedure approach? Some positives and negatives. ICES Journal of Marine Science, 64: 613--617.

Butterworth, D. S., and Punt, A. E. 1999. Experiences in the evaluation and implementation of management procedures. ICES Journal of Marine Science, 56: 985--998.

Cadigan, N. G. 2016. A state-space stock assessment model for northern
cod, including under-reported catches and variable natural mortality
rates. Canadian Journal of Fisheries and Aquatic Sciences, 73: 296--308.

Ch\'avez, F. P., Ryan, J., Lluch-Cota, S. E., and Niquen, M. 2003. From
anchovies to sardines and back: multidecadal change in the Pacific
Ocean. Science, 299: 217--221.

Chen, T., and Francis, B. A. 1995. Optimal Sampled-Data Control Systems. Springer, London.

Cohen, J. 1988. Statistical Power Analysis for the Behavioral Sciences,
2nd edn. Lawrence Erlbaum, Hillsdale, NJ.

Costantino, R. F., Cushing, J. M., Dennis, B., and Desharnais, R. A.
1995. Experimentally induced transitions in the dynamic behaviour of
insect populations. Nature, 375: 227--230.

DFO. 2011. Recovery potential assessment for the Newfoundland and
Labrador Designatable Unit (NAFO Divs. 2GHJ, 3KLNO) of Atlantic Cod. DFO
Canadian Science Advisory Secretariat Science Advisory Report 2011/037.

DFO. 2016. Stock Assessment of Northern cod (NAFO Divs. 2J3KL) in 2016.
DFO Canadian Science Advisory Secretariat Science Advisory Report
2016/026.

DFO. 2022. Stock Assessment of Northern cod (NAFO Divs. 2J3KL) in 2022.
DFO Canadian Science Advisory Secretariat Science Advisory Report
2022/041.

DFO. 2024. NAFO Divisions 2J3KL Northern cod (Gadus morhua) stock
assessment to 2024. DFO Canadian Science Advisory Secretariat Science
Advisory Report.

Forssell, U., and Ljung, L. 1999. Closed-loop identification revisited.
Automatica, 35: 1215--1241.

Gurney, W. S. C., Blythe, S. P., and Nisbet, R. M. 1980. Nicholson's
blowflies revisited. Nature, 287: 17--21.

Guti\'errez, D., Sifeddine, A., Field, D. B., Ortlieb, L., Vargas, G.,
Ch\'avez, F. P., Velazco, F., Ferreira, V., Tapia, P., Salvatteci, R.,
Boucher, H., Morales, M. C., Vald\'es, J., Reyss, J.-L., Campusano, A.,
Boussafir, M., Mandeng-Yogo, M., Garc\'ia, M., and Baumgartner, T. 2009.
Rapid reorganization in ocean biogeochemistry off Peru towards the end
of the Little Ice Age. Biogeosciences, 6: 835--848.

ICES. 2020. Inter-Benchmark Process on BAltic Sprat (Sprattus sprattus)
and Herring (Clupea harengus) (IBPBash). ICES Scientific Reports. 2:34.
44 pp.

ICES. 2021. Arctic Fisheries Working Group (AFWG). ICES Scientific
Reports. 3:58. 817 pp.

Karlsson, M., and Gilek, M. 2020. Mind the gap: Coping with delay in environmental governance. Ambio, 49: 1067--1075.

Lomb, N. R. 1976. Least-squares frequency analysis of unequally spaced
data. Astrophysics and Space Science, 39: 447--462.

Ludwig, D., Jones, D. D., and Holling, C. S. 1978. Qualitative analysis
of insect outbreak systems: the spruce budworm and forest. Journal of
Animal Ecology, 47: 315--332.

McManus, M. C., Licandro, P., and Coombs, S. H. 2016. Is the Russell
Cycle a true cycle? Multidecadal zooplankton and climate trends in the
western English Channel. ICES Journal of Marine Science, 73: 227--238.

Moxnes, E. 1998. Not only the tragedy of the commons: misperceptions of
bioeconomics. Management Science, 44: 1234--1248.

Ne\v{s}i\'c, D., and Teel, A. R. 2004. A framework for stabilisation of
nonlinear sampled-data systems based on their approximate discrete-time
models. IEEE Transactions on Automatic Control, 49: 1103--1122.

Pauly, D., and Tsukayama, I. (eds) 1987. The Peruvian anchoveta and its
upwelling ecosystem: three decades of change. ICLARM Studies and
Reviews, 15.

Punt, A. E., and Donovan, G. P. 2007. Developing management procedures
that are robust to uncertainty: lessons from the International Whaling
Commission. ICES Journal of Marine Science, 64: 603--612.

Punt, A. E., Butterworth, D. S., de Moor, C. L., De Oliveira, J. A. A., and Haddon, M. 2016. Management strategy evaluation: best practices. Fish and Fisheries, 17: 303--334.

Ricard, D., Minto, C., Jensen, O. P., and Baum, J. K. 2012. Examining
the knowledge base and status of commercially exploited marine species
with the RAM Legacy Stock Assessment Database. Fish and Fisheries, 13:
380--398.

Rose, G. A. 2026. Northern cod comeback: 10 years after. Canadian
Journal of Fisheries and Aquatic Sciences, 83: 1--14.
https://doi.org/10.1139/cjfas-2025-0141

Rose, G. A., and Rowe, S. 2015. Northern cod comeback. Canadian Journal
of Fisheries and Aquatic Sciences, 72: 1789--1798.

Rose, G. A., and Walters, C. J. 2019. The state of Canada's iconic
Northern cod: a second opinion. Fisheries Research, 219: 105314.

Scargle, J. D. 1982. Studies in astronomical time series analysis. II.
Statistical aspects of spectral analysis of unevenly spaced data. The
Astrophysical Journal, 263: 835--853.

Statistics Canada. Tables 38-10-0167-01 and 38-10-0168-01, CANSIM
database. Statistics Canada, Ottawa.

Shertzer, K. W., and Prager, M. H. 2007. Delay in fishery management: diminished yield, longer rebuilding, and increased probability of stock collapse. ICES Journal of Marine Science, 64: 149--159.

Stuart, A. M., and Humphries, A. R. 1996. Dynamical Systems and Numerical Analysis. Cambridge University Press, Cambridge.

Sumaila, U. R., Ebrahim, N., Schuhbauer, A., Skerritt, D., Li, Y., Kim, H. S., Mallory, T. G., Lam, V. W. L., and Pauly, D. 2019. Updated estimates and analysis of global fisheries subsidies. Marine Policy, 109: 103695.

Tam, J. C., and Bundy, A. 2019. Mass-balance models of the Newfoundland
and Labrador Shelf ecosystem for 1985--1987 and 2013--2015. Canadian
Technical Report of Fisheries and Aquatic Sciences, 3328.

UNFCCC. 2015. Paris Agreement. United Nations Framework Convention on Climate Change.

Walters, C. J. 1986. Adaptive Management of Renewable Resources. Macmillan, New York.

World Bank. Poverty and Inequality Platform. World Bank, Washington, DC.


\section*{Declarations}
\begin{description}[leftmargin=2.0cm,style=nextline]
\item[Funding] No external funding was received for this work.
\item[Conflicts of interest] The author declares no competing interests.
\item[Data and code availability] The source bundle contains the model core, analysis scripts, archived tables, recovered local inputs used for re-execution, and the Supplementary Information. The RAM Legacy database is identified by version and source; upstream data remain subject to their original terms. A reconstruction register distinguishes newly reproduced outputs from archived records and unresolved inputs. The full 17-cell white-noise power table is assembled from the archived 15-cell run log and the two-cell cod follow-up JSON; a complete rerun of the driver did not finish. This limitation is stated in Supplementary S13.1.
\item[Author contributions] A.A. conceived the analysis, developed the model, conducted the re-analysis, and wrote the manuscript.
\item[Ethics approval] Not applicable.
\end{description}
\end{document}
